\documentclass{article}
\usepackage{amsmath}
\usepackage{amssymb}
\usepackage{graphicx}
\usepackage{hyperref}
\usepackage{geometry}
\geometry{margin=1in}

\begin{document}

\title{Problem Set 2}
\author{Tate Mason}
\date{\today}
\maketitle

\section*{Problem 1: Data Overview and Analysis}

\subsection*{Part A: Data Summary}

In Colombia, we see fewer overall observations in every sector:
\begin{itemize}
  \item Food Industry: 6140
  \item Apparel Industry: 4454
  \item Metal Industry: 3678 
  \item Textile Industry: 2840 
  \item Wood Industry: 964
\end{itemize}

So, we see that Colombia has a huge food and manufacturing sector, whereas they have a smaller industry producing raw materials like wood.

In Chile, we see an interesting distribution of observations:
\begin{itemize}
  \item Food Industry: 18833
  \item Metal Industry: 5856
  \item Textile Industry: 5447
  \item Wood Industry: 4835 
  \item Apparel Industry: 4694 
\end{itemize}

So, Chile has a massive food industry, but also has a balanced manufacturing sector with a lot of metal, textile, and wood production.

Over time, we see varying trends, indicating some industrial growth and some decline in different sectors for both countries. Further, what may be a data aggregation
issue, 1981 has disproportionately more observations than other years for both countries. Something to look back at.

\begin{tabular}{l|c|c}
  \textbf{Year} & \textbf{Colombia Observations} & \textbf{Chile Observations} \\
  \hline
  1979 & 

\end{tabular}

\subsection*{Part B: Analysis of Investment by Industry}

\includegraphics[width=\linewidth]{~/SchoolWork/Y2S1/IO/ps2/1.2_col.pdf}
  \hfill
\includegraphics[width=\linewidth]{~/SchoolWork/Y2S1/IO/ps2/1.2_chi.pdf}

In Colombia, the is much more dispersion in investment levels across industries, with someindustries even experiencing negative investment. Food, metal,
and textiles have very similarinvestment patterns, with means around 5. Apparel and wood are similar to the other three,though a bit lower. In Chile,
investment levels are super consistent across industries, with allindustries having means aound 5-7. Distributions are also much tighter, though most do
have some negative investment observations.

\subsection*{Part C: Analysis of Ouput by Industry}

\includegraphics[width=\linewidth]{~/SchoolWork/Y2S1/IO/ps2/1.3_col.pdf}
  \hfill
\includegraphics[width=\linewidth]{~/SchoolWork/Y2S1/IO/ps2/1.3_chi.pdf}

In Colombia, output levels do not vary too much across industries, with means around 7-10and similar distributions. In Chile, distribtions are also similar,
though there is a bit morevariation. Textiles and food have a wider distribution of output levels, while apparel and metal are a bit tighter.

\subsection*{Part D: Analysis of Labor Output Correlation by Industry}

\includegraphics[width=\linewidth]{~/SchoolWork/Y2S1/IO/ps2/1.4_col.pdf}
  \hfill
\includegraphics[width=\linewidth]{~/SchoolWork/Y2S1/IO/ps2/1.4_chi.pdf}

In both countries, there is a positive relationship between labor input and output across allindustries. The strength of this relationship varies by industry,
with some industries showing astronger correlation than others. For example, in Colombia, the food industry shows a strongpositive correlation, while the wood
industry has a weaker correlation. In Chile, all industries show a strong positive correlation.

\section*{Problem 2: Estimation Techniques}
\subsection*{Part A: OLS}

To estimate via OLS, we estimate the following equation:
\begin{equation}
  Y_{it} = e^{\alpha_0 + \omega_{it} + \varepsilon_{it}} K_{it}^{\alpha_k} L_{it}^{\alpha_l}
\end{equation}

\begin{equation}
  y_{it} = \alpha_0 + \omega_{it} + \alpha_k k_{it} + \alpha_l l_{it} + \varepsilon_{it}
\end{equation}
Where $y$ is log of value-added output, $k$ is log capital, $l$ is log labor, and $\omega$ is the productivity shock. We estimate this equation for both Colombia and Chile,
first homogenous across industries, then heterogenous by industry.

\textbf{OLS Homogenous Results:}

\begin{center}
\begin{tabular}[h]{|lcccc|}
  \hline
  \textbf{Country} & \textbf{$\alpha_0$} & \textbf{$\alpha_k$} & \textbf{$\alpha_l$} & \textbf{R$^2$} \\
  \hline
  Colombia (Homogenous) & 2.186 & 0.273 & 0.831 & 0.823 \\
  Chile (Homogenous) & 2.351 & 0.887 & 0.303 & 0.660 \\
  \hline
\end{tabular}
\end{center}

\vspace{1cm}
\textbf{OLS Heterogenous Results:}
\begin{center}
\begin{tabular}[h]{|lccccc|}
  \hline
  \textbf{Country} & \textbf{Industry} & \textbf{$\alpha_0$} & \textbf{$\alpha_k$} & \textbf{$\alpha_l$} & \textbf{R$^2$} \\
  \hline
  \hline
  Colombia & & & & & \\
  \hline
   & Food & 2.238 & 0.321 & 0.754 & 0.831 \\
   & Apparel & 2.510 & 0.133 & 0.921 & 0.831 \\
   & Metal & 1.974 & 0.208 & 0.998 & 0.831 \\
   & Textile & 2.316 & 0.257 & 0.858 & 0.831 \\
   & Wood & 2.559 & 0.117 & 0.858 & 0.831 \\
   \hline
  Chile & & & & & \\
  \hline
   & Food & 1.950 & 0.334 & 0.896 & 0.689 \\
   & Apparel & 2.865 & 0.150 & 1.055 & 0.689 \\
   & Metal & 2.734 & 0.232 & 0.994 & 0.689 \\
   & Textile & 2.709 & 0.243 & 0.874 & 0.689 \\
   & Wood & 2.488 & 0.174 & 1.016 & 0.689 \\
   \hline
\end{tabular}
\end{center}

These results indicate that there are significant differences in the production function parameters across countries and industries. Colombia has a higher labor
elasticity compared to Chile, suggesting that labor plays a more important role in production in Colombia. Conversely, Chile has a higher capital elasticity, 
indicating that capital is more important in their production process. When looking at heterogenous estimates, we see that within each country, different industries 
have varying elasticities, reflecting the unique production processes and input requirements of each industry.

\subsection*{Part B: OLS with Investment}

To estimate via OLS with investment, we modify our previous equation to include investment as an additional regressor:
\begin{equation}
  y_{it} = \alpha_0 + \omega_{it} + \alpha_k k_{it} + \alpha_l l_{it} + \alpha_i i_{it} + \varepsilon_{it}
\end{equation}
Where $i$ is log investment. We estimate this equation for both Colombia and Chile, first homogenous across industries, then heterogenous by industry.

\textbf{OLS with Investment Homogenous Results:}
\begin{center}
\begin{tabular}[h]{|lccccc|}
  \hline
  \textbf{Country} & \textbf{$\alpha_0$} & \textbf{$\alpha_k$} & \textbf{$\alpha_l$} & \textbf{$\alpha_i$} & \textbf{R$^2$} \\
  \hline
  Colombia (Homogenous) & 2.282 & 0.242 & 0.828 & 0.031 & 0.836 \\
  Chile (Homogenous) & 2.521 & 0.299 & 0.844 & 0.017 & 0.689 \\
  \hline
\end{tabular}
\end{center}

\vspace{1cm}
\textbf{OLS with Investment Heterogenous Results:}
\begin{center}
\begin{tabular}[h]{|lcccccc|}
  \hline
  \textbf{Country} & \textbf{Industry} & \textbf{$\alpha_0$} & \textbf{$\alpha_k$} & \textbf{$\alpha_l$} & \textbf{$\alpha_i$} & \textbf{R$^2$} \\
  \hline
  \hline
  Colombia & & & & & & \\
  \hline
   & Food & 2.384 & 0.267 & 0.767 & 0.042 & 0.844 \\
   & Apparel & 2.545 & 0.119 & 0.925 & 0.015 & 0.844 \\
   & Metal & 2.041 & 0.196 & 0.985 & 0.017 & 0.844 \\
   & Textile & 2.278 & 0.230 & 0.834 & 0.045 & 0.844 \\
   & Wood & 2.518 & 0.139 & 0.954 & -0.001 & 0.844 \\
   \hline
  Chile & & & & & & \\
  \hline
   & Food & 2.021 & 0.300 & 0.882 & 0.020 & 0.702 \\
   & Apparel & 2.925 & 0.120 & 1.032 & 0.018 & 0.702 \\
   & Metal & 2.784 & 0.200 & 0.980 & 0.019 & 0.702 \\
   & Textile & 2.759 & 0.210 & 0.860 & 0.021 & 0.702 \\
   & Wood & 2.538 & 0.150 & 1.002 & 0.022 & 0.702 \\
   \hline
\end{tabular}
\end{center}

\subsection*{Part C: OP Estimation}

\subsection*{Part D: LP Estimation}

\subsection*{Part E: ACF Estimation}

\section*{Problem 3: Productivity Dynamics}

\subsection*{Part A: Estimation of $\omega$ Using OP Estimates - Homogenous}

\begin{center}
\begin{tabular}[h]{|lcc|}
  \hline
  \textbf{Country} & \textbf{Mean of $\omega$} & \textbf{Std. Dev. of $\omega$} \\
  \hline
  Colombia & 294.294 & 702.900 \\
  Chile & 74.876 & 82.639 \\
  \hline
\end{tabular}
\end{center}

In Colombia, we see a much higher mean productivity level compared to Chile, with a much larger standard deviation as well. This indicates that Colombian firms
are generally more productive, but also have a wider range of productivity levels compared to Chilean firms. Thus, Colombia's calculation seems to have some sort
of outlier issue, as the standard deviation is more than double the mean. This warrants further investigation into the data and estimation process to ensure accuracy.

\subsection*{Part B: Estimation of $\omega$ Using OP Estimates - Heterogenous}

\begin{center}
\begin{tabular}[h]{|lccc|}
  \hline
  \textbf{Country} & \textbf{Industry} & \textbf{Mean of $\omega$} & \textbf{Std. Dev. of $\omega$} \\
  \hline
  Colombia & & & \\
  \hline
   & Food & 61.7 & 115 \\
   & Apparel & 163 & 376 \\
   & Metal & 238 & 400 \\
   & Textile & 124 & 279 \\
   & Wood & 294 & 1310 \\
  Chile & & & \\
  \hline
   & Food & 90.9 & 124 \\
   & Apparel & 460 & 1849 \\
   & Metal & 88.7 & 119 \\
   & Textile & 62.7 & 134 \\
   & Wood & 357 & 619 \\
    \hline
\end{tabular}
\end{center}

These are super weird results due to the massive standard deviations, likely caused by outliers in the data. What can be seen is that Colombia has a generally
higher mean productivity than Chile across all industries except wood and apparel. Within Colombia, wood and metal industries have the highest mean productivity,
while food has the lowest. In Chile, apparel and wood have the highest mean productivity, while textiles and metal have the lowest. 

\subsection*{Part C: Time Trends in Productivity}

\subsubsection*{Homogenous Productivity Trends}

\includegraphics[width=\linewidth]{~/SchoolWork/Y2S1/IO/ps2/3.3_hom_col.pdf}
\includegraphics[width=\linewidth]{~/SchoolWork/Y2S1/IO/ps2/3.3_hom_chi.pdf}

\subsubsection*{Heterogenous Productivity Trends by Industry}

\includegraphics[width=\linewidth]{~/SchoolWork/Y2S1/IO/ps2/3.3_ind_col_mean.pdf}
\includegraphics[width=\linewidth]{~/SchoolWork/Y2S1/IO/ps2/3.3_ind_col_sd.pdf}

\includegraphics[width=\linewidth]{~/SchoolWork/Y2S1/IO/ps2/3.3_ind_chi_mean.pdf}
\includegraphics[width=\linewidth]{~/SchoolWork/Y2S1/IO/ps2/3.3_ind_chi_sd.pdf}



\subsubsection*{Discussion}

In Colombia, average productivity trends upward over the time horizon, with a slightly lesser increase in standard deviation, indicating that the firms
are becoming more productive on average with time, while variability is not increasing to the same degree. In Chile, we see a a spike in productivity after
1985, continuing until 1990 before slowing. Standard deviation also increases over time, with more consistency in growth. This suggests that Chilean
firms are also becoming more productive on average, but with greater variability among firms.

When splitting out by industry, we see that Colombian textile industries are the most productive on average, while wood and apparel are the least. However,
teh textile industry also exhibits the highest variability, while food is the least so. In Chile, all industries are increasing in productivity over time to similar
degrees, with textiles being the most productive until about 1992 when food overtakes it. Apparel is the least productive until about 1994, when metal
takes over as least prductive. Food and apparel are the most variable industries in Chile when looking at growth in standard deviation, though textiles
have the highest level until about 1990. Over the whole horizon, wood and metal are the least variable.

\end{document}
